\section{Modes that depend on the deflating coordinates}
\label{sec:modes}

\subsection{Three kinds of modes}

Modes fall into three kinds. A mode that does not depend on the deflating coordinates,
$k_5=k_6=k_7=0$, obeys the flat equation; its solutions are the plane waves of the earlier paper, with
the energies $\pm\Eperp$, eight of each sign. A mode with some $k_t\ne0$ behaves in the same way as
long as its physical wavenumber is negligible, $\kd\ll\Eperp$. But $\kd$ grows without bound, and
every such mode sooner or later reaches $\kd\gg\Eperp$. The moment at which $\kd=\Eperp$ is its
\emph{crossing}. \SIG{It is the turning value (\ref{eq:turn}) of the classical particle, with the
wavenumbers in place of the momenta.}{} The crossing is where the deflation acts on the field, and
this section solves the equation of the weighted field across it.

The solution is exact when \emph{one rate governs}: when only one of the $k_t$ differs from zero, or
when the three rates are equal. In both cases the matrix \SIG{$\gamma^{d}$}{$\alpha^{d}$} does not
depend on the time, and $\kd$ is proportional to $e^{\rt\,x4}$ with the rate that governs. The natural
variable is then the physical wavenumber in units of the rate,
\begin{equation}
\ud=\frac{\kd(x4)}{\rt},\qquad\frac{d}{dx4}=\rt\,\ud\,\frac{d}{d\ud}.
\label{eq:ud}
\end{equation}
The crossing is at $\ud=\Eperp/\rt$, and this ratio of the energy to the rate is the one parameter
on which the results depend.

%%<1671>
\subsection{The exact solution}

In $\hh=\hh_\perp+\kd\,\alpha^{d}$ the two constant matrices satisfy $\hh_\perp^{2}=\Eperp^{2}$ and
$(\alpha^{d})^{2}=1$, and they anticommute. Let $w$ be one of the eight eigenvectors of $\hh_\perp$
with the eigenvalue $+\Eperp$; then $\alpha^{d}w$ is an eigenvector with the eigenvalue $-\Eperp$.
Let $Z_+$ and $Z_-$ denote two cylinder functions of the same kind (Bessel functions of the first or
of the second kind, or Hankel functions) of the orders $\tfrac12+i\Eperp/\rt$ and
$-\tfrac12+i\Eperp/\rt$. The recurrence relations of the cylinder functions connect two functions
whose orders differ by one, and with them one verifies by substitution that
\begin{equation}
\hpsi=\sqrt{\ud}\,\Big[Z_+(\ud)\,w+i\,Z_-(\ud)\,\alpha^{d}w\Big]
\label{eq:exact}
\end{equation}
solves Eq.~(\ref{eq:h}) for each of the eight $w$. Two independent kinds of cylinder function give
all sixteen solutions. Equation (\ref{eq:exact}) was verified with the Bessel functions of an
arbitrary-precision library, to a relative residual of $10^{-15}$.

\subsection{The deflation creates particles}

At early times, $\ud\to0$, a Bessel function of the first kind behaves as the power of its argument
that its order gives. The solution that has positive energy at early times,
$\propto e^{-i\Eperp x4}\,w$, is therefore Eq.~(\ref{eq:exact}) with $Z_+$ the Bessel function of the
first kind of the order $-\tfrac12-i\Eperp/\rt$ and $Z_-$ minus that of the order
$+\tfrac12-i\Eperp/\rt$, a pair that obeys the same recurrence relations. At late times, $\ud\to\infty$, the solution of
positive energy is Eq.~(\ref{eq:exact}) with the Hankel functions of the second kind, which behave as
$e^{-i\ud}$, and the solution of negative energy is the one with the Hankel functions of the first
kind. The connection formula of the Bessel functions expresses the first solution through the other
two, and the weight of the negative-energy part divided by that of the positive-energy part is
$e^{-2\pi\Eperp/\rt}$. A mode that starts empty therefore ends with
\begin{equation}
\ncr=\frac{1}{e^{2\pi\Eperp/\rt}+1}
\label{eq:ncr}
\end{equation}
particles in each of its eight internal states, and with as many antiparticles. This is the
Fermi-Dirac distribution in the energy $\Eperp$ at the temperature
\begin{equation}
\Tdefl=\frac{\rt}{2\pi}.
\label{eq:Tdefl}
\end{equation}
For equal rates it is the temperature that Gibbons and Hawking found for de Sitter space
\cite{GibbonsHawking1977}; the creation of particles by a time-dependent metric is the phenomenon
discovered by Parker \cite{Parker1969,BirrellDavies1982}.

\begin{figure*}
\includegraphics[width=\textwidth]{fig_creation}
\caption{\label{fig:creation}Creation of particles on the deflating $7+1$ spacetime. (a) Created
quanta per state, multiplied by $e^{2\pi\,\mathrm{ratio}}$, against the ratio of the energy to the
rate. Fermions follow the Fermi-Dirac law (\ref{eq:ncr}), the modes of the Vorton-Scalar the
Bose-Einstein law (\ref{eq:bose}), both at the temperature $\gff/2\pi$. Lines: the exact laws;
symbols: computed. (b) History of two modes, with $\Eperp/\gff=1$ and $0.5$: the number of particles
in the mode against the physical wavenumber $\kd/\Eperp$, which grows with the time. After the
crossing at $\kd=\Eperp$ the number settles at the value of Eq.~(\ref{eq:ncr}) (dotted). (c) Two
deflating directions, $\gff=1$ and $\gss=2$, $1.2$ or $0.5$: mean number of created particles per
state against the delay between the two crossings; the dotted lines are Eq.~(\ref{eq:ncr}) with each
of the two rates.}
\end{figure*}

The numerical solution of Eq.~(\ref{eq:h}) for all sixteen components confirms the law
(Fig.~\ref{fig:creation}). For $\Eperp/\gff$ between 0.25 and 3 the created number agrees with
Eq.~(\ref{eq:ncr}) to a relative deviation of at most $3.2\times10^{-6}$, while the numbers
themselves range from 0.17 to $6.5\times10^{-9}$; the eight internal states are created equally.
With three equal rates the result depends neither on the direction of the wavenumber nor on the
constants $\ct$. With unequal rates no exact solution is known. If the second physical wavenumber is
still below $e^{-3}\Eperp$ when the first reaches $\Eperp$, the mean created number is
Eq.~(\ref{eq:ncr}) with the rate of the coordinate that becomes energetic first, to within 1.2 per
cent. If two directions become energetic together, the eight internal states are created
\emph{unequally}, in the cases computed by up to a factor 110 between the most and the least
populated state, and the result is known only numerically.

\subsection{What the result means}

Each mode with $k_t\ne0$ is populated once, at its crossing. Its particles then gain energy without
bound, as $(\Eperp^{2}+\kd^{2})^{1/2}$, like the classical particle of Sec.~\ref{sec:geodesics}.
Modes with $k_t=0$ stay empty. Two remarks are in order.

First, the stress of the created particles must eventually change the metric. That reaction of the
fields on the metric has not been computed, and the results of this section hold for as long as it
can be neglected. For any rate of deflation that the bounds of Sec.~\ref{sec:cosmo} allow today the
temperature (\ref{eq:Tdefl}) is below $10^{-30}$~K, and nothing is created.

Second, the instantaneous energy of a quantum in a mode with $k_t\ne0$ is at least $|k_t|/b_t$,
which grows as $e^{\rt\,x4}$. A process with a given energy at its disposal can excite such a mode
only before the mode has crossed that energy. If the deflating coordinates are periodic, the
smallest wavenumber that differs from zero is fixed, and the sector of fields that do not depend on
$x5,x6,x7$ is separated from the rest by a gap that grows exponentially. This is the sense in which
the deflation supports the restriction to that sector, which the earlier paper assumed. If the
deflating coordinates are infinite, modes of arbitrarily small $k_t$ exist at every time, and the
restriction remains an assumption.
%%</1671>
%%<1644>
\subsection{The exact solution}

In $K=K_\perp+i\,\kd\,\gamma^{4}\gamma^{d}$ the first matrix is anti-Hermitian with
$K_\perp^{2}=-\Eperp^{2}$, the matrix $i\gamma^{4}\gamma^{d}$ is Hermitian with the square one, and
the two anticommute. Let $w$ be one of the eight eigenvectors of $K_\perp$ with the eigenvalue
$+i\Eperp$; then $i\gamma^{4}\gamma^{d}w$ is an eigenvector with the eigenvalue $-i\Eperp$ and is
orthogonal to $w$. The equation is solved by the modified Bessel functions that decay at large
argument, of the orders $\tfrac12\mp i\Eperp/\rt$. Such a function is written here with its order
as a subscript and with its argument; the letter $K$ without subscript and argument is the matrix of
Eq.~(\ref{eq:K}). With the recurrence relations of these functions one verifies by substitution that
\begin{equation}
\begin{aligned}
\hpsi=\sqrt{\ud}\,\Big[&K_{\frac12-i\Eperp/\rt}(\ud)\,w\\
&-K_{\frac12+i\Eperp/\rt}(\ud)\,i\gamma^{4}\gamma^{d}w\Big]
\end{aligned}
\label{eq:exact}
\end{equation}
solves Eq.~(\ref{eq:K}) for each of the eight $w$; the relative residual, computed with the Bessel
functions of an arbitrary-precision library, is $2\times10^{-15}$. The other eight solutions are built in
the same way from the modified Bessel functions that \emph{grow} as $e^{+\ud}$.

\subsection{Standing waves, and their extinction}

Since $\ud$ itself grows exponentially with the time, a solution of the second kind grows as the
exponential of an exponential. Such a solution is not admissible; it is the counterpart of the
growing exponential under a potential barrier. Admissible are the eight solutions
(\ref{eq:exact}), which stay bounded. Their properties are exact.

\emph{Magnitude.} Because $w$ and $i\gamma^{4}\gamma^{d}w$ are orthogonal, and because the two
functions in Eq.~(\ref{eq:exact}) are complex conjugates of each other,
\begin{equation}
|\hpsi|^{2}=2\,\ud\,\big|K_{\frac12+i\Eperp/\rt}(\ud)\big|^{2}\,|w|^{2}.
\label{eq:magnitude}
\end{equation}

\emph{Early times}, $\ud\to0$: $|\hpsi|^{2}\to\pi\,|w|^{2}/\cosh(\pi\Eperp/\rt)$, a constant. The
solution is then a combination of $e^{+i\Eperp x4}\,w$ and $e^{-i\Eperp x4}\,i\gamma^{4}\gamma^{d}w$
with coefficients of \emph{equal} magnitude: the two signs of the frequency enter with equal weight.
A bounded solution is a \emph{standing wave} in the ordinary time, not a traveling wave. Its
conserved charge $\hpsi^\dagger W_{1644}\hpsi$ vanishes.

\emph{Late times}, $\ud\to\infty$: $|\hpsi|^{2}\to\pi\,|w|^{2}e^{-2\ud}$. The mode dies out as the
exponential of an exponential of the time.

\begin{figure*}
\includegraphics[width=\textwidth]{fig_extinction}
\caption{\label{fig:extinction}Extinction on the deflating $4+4$ spacetime. (a) A bounded solution:
$|\hpsi|^{2}$ against the physical wavenumber $\kd/\Eperp$, which grows with the time; computed
(line) and the exact form (\ref{eq:magnitude}) (symbols). It is constant long before the crossing at
$\kd=\Eperp$ and falls after it. (b) Five cases, with one, two or three wavenumbers $k_t$ different
from zero and with equal and unequal rates, on a logarithmic scale. (c) The bounded modes of the
Vorton-Scalar on the same spacetime (Sec.~\ref{sec:scalarmodes}), which are extinguished in the same
way; lines: exact, symbols: computed.}
\end{figure*}

\emph{The law of the extinction.} In the semiclassical (WKB) approximation the bounded solution
falls beyond the turning value, at which $\kd=\Eperp$, as $e^{-\Swkb}$ with
\begin{align}
\Swkb&=\int_{x4_{\mathrm{turn}}}^{x4}\sqrt{\kd^{2}-\Eperp^{2}}\;dx4'\nonumber\\
&=\frac{1}{\rt}\Big[\sqrt{\kd^{2}-\Eperp^{2}}-\Eperp\arccos\frac{\Eperp}{\kd}\Big],
\label{eq:Swkb}
\end{align}
which approaches $(\kd-\tfrac{\pi}{2}\Eperp)/\rt$ for $\kd\gg\Eperp$. The exact solution fixes the
constant that the approximation leaves open:
\begin{equation}
\begin{aligned}
\ln\frac{|\hpsi|_{\mathrm{early}}}{|\hpsi|_{\mathrm{late}}}-\Swkb\ \longrightarrow\ {}&\tfrac12\ln2\\
&-\tfrac12\ln\big(1+e^{-2\pi\Eperp/\rt}\big).
\end{aligned}
\label{eq:constant}
\end{equation}

The numerical solution of Eq.~(\ref{eq:K}) for all sixteen components confirms these statements
(Fig.~\ref{fig:extinction}). Beyond the turning value exactly eight of the sixteen solutions stay
bounded; followed backwards to early times they contain parts of positive and of negative frequency
of equal weight, to $2\times10^{-14}$. Where one rate governs, the computed magnitude equals
Eq.~(\ref{eq:magnitude}) over the range $0.02<\ud<12$ to a relative deviation of
$3.6\times10^{-12}$, and Eq.~(\ref{eq:constant}) holds to $3.6\times10^{-4}$ at $\Swkb=30$. With
three unequal rates no exact solution is known; the bounded solutions fall as $e^{-\Swkb}$ times a
factor that depends on the internal state, with a spread of 0.40 in its logarithm at $\Swkb=30$.

\subsection{What the result means}

On the flat $4+4$ spacetime a field that depends on the additional times evolves non-unitarily in
the ordinary time. This is the principal difficulty of four time-like dimensions
\cite{Tegmark1997,CraigWeinstein2009}, and the earlier paper met it by imposing that the fields do
not depend on those times \cite{Flat1644}. On the deflating $4+4$ spacetime every admissible mode
that depends on a deflating time is a standing wave that is extinguished after its crossing, with an
exponent $\Swkb$ that grows as $e^{\rt\,x4}$. What remains at late times is the sector of fields
that do not depend on $x5,x6,x7$, and in that sector the evolution is unitary and
$\hpsi^\dagger\hpsi$ is conserved. \emph{The deflation of the three additional times removes,
dynamically, the sector that causes the difficulty.} The classical counterpart is the turning value
of Sec.~\ref{sec:geodesics}.

The limits of this conclusion must be stated with equal clarity. The eight growing solutions are
discarded because they are unbounded. That is a condition of admissibility, imposed as in the theory
of a potential barrier; it is not derived from a principle. The theory still has no positive
conserved charge for modes that depend on the deflating times, and a consistent quantization exists,
as before, only for the sector that does not depend on them. What has been shown is that the
admissible modes outside that sector die out, not that the difficulty of four time-like dimensions
has been solved.
%%</1644>

\subsection{The numerical method}

The equations of this section are linear, with coefficients that grow exponentially, and an
integrator with a fixed step fails on them. The program uses a commutator-free Magnus method of the
fourth order: over each step the solution is advanced by the product of two matrix exponentials of
linear combinations of the coefficient matrix at two Gauss points. Because the square of
\SIG{$K$}{$\hh$} is a multiple of the unit matrix at any one time, each exponential is evaluated in
closed form\SIG{.}{, and every step is exactly unitary.} The length of a step is chosen so that
neither the phase accumulated in it nor the relative change of the coefficients exceeds a fixed
bound. \SIG{}{The number of particles is read off by projection on the eigenstates of the
Hamiltonian corrected to the first order in the rate at which its direction changes, which removes
the leading spurious oscillation of the particle number at finite time. }An independent program,
which uses an explicit Runge-Kutta method of the eighth order and the Bessel functions of another
library, \SIG{confirms Eq.~(\ref{eq:exact}) and, to the accuracy of the semiclassical form at the
exponent that it reaches, Eq.~(\ref{eq:constant})}{reproduces Eq.~(\ref{eq:ncr}) to
$2.2\times10^{-5}$} \cite{DeflationRepo}.

